%% =====================================================================
%%  GENERATED FILE — do not edit by hand.
%%  Source: data/profile.yaml   Template: templates/cv.tex.j2
%%  Rebuild: python scripts/build.py && latexmk -pdf cv/Samuel_Adeniji_CV.tex
%% =====================================================================
\documentclass[11pt,a4paper]{article}

\usepackage{iftex}
\ifPDFTeX
  \usepackage[T1]{fontenc}
  \usepackage[utf8]{inputenc}
\else
  \usepackage{fontspec}
\fi
\usepackage[margin=1.9cm]{geometry}
\usepackage{titlesec}
\usepackage{enumitem}
\usepackage{xcolor}
\usepackage{tabularx}
\usepackage[hidelinks]{hyperref}

\definecolor{accent}{HTML}{1F3A5F}
\hypersetup{colorlinks=true, urlcolor=accent, linkcolor=accent}

\titleformat{\section}
  {\normalfont\large\bfseries\color{accent}}
  {}{0em}{\MakeUppercase}[\vspace{-0.7em}\rule{\textwidth}{0.8pt}]
\titlespacing*{\section}{0pt}{1.1em}{0.6em}

\newcommand{\entry}[2]{%
  \noindent\begin{tabularx}{\textwidth}{@{}X r@{}}
  \textbf{#1} & \textit{#2}
  \end{tabularx}\par\vspace{0.15em}}

\newcommand{\subentry}[1]{\noindent\textit{\small #1}\par\vspace{0.2em}}

\newcommand{\stack}[1]{\noindent{\small\textbf{\color{accent}Stack:} #1}\par\vspace{0.15em}}

\setlist[itemize]{leftmargin=1.2em, itemsep=0.12em, topsep=0.2em, parsep=0pt}
\newenvironment{bullets}{\begin{itemize}\small}{\end{itemize}\vspace{0.2em}}

\pagestyle{empty}
\setlength{\parindent}{0pt}


\begin{document}

\begin{center}
  {\LARGE\bfseries\color{accent} SAMUEL AKINWUMI ADENIJI}\\[0.35em]
  {\small AI for Healthcare $\cdot$ Multimodal Medical AI $\cdot$ Medical Image Analysis}\\[0.35em]
  {\small
    \href{mailto:sadeniji@andrew.cmu.edu}{sadeniji@andrew.cmu.edu} $\mid$
    Carnegie Mellon University $\mid$
    \href{https://sadeniji.vercel.app/}{Website} $\cdot$
    \href{https://scholar.google.com/citations?hl=en&user=oM69cJoAAAAJ}{Google Scholar} $\cdot$
    \href{https://www.linkedin.com/in/saadeniji/}{LinkedIn} $\cdot$
    \href{https://github.com/Tech-sam-90}{GitHub}
  }
\end{center}
\vspace{0.4em}

\section{Research Interests}
Medical vision--language models for 3D imaging; instruction-conditioned and token-efficient architectures for volumetric data; grounded radiology report generation and multi-scale representation learning; interpretable and reliable medical AI; deployment of imaging AI in low-resource and African clinical settings.

\section{Education}

\entry{Carnegie Mellon University --- Master of Science, Engineering Artificial Intelligence}{Aug 2025 -- May 2027}
\subentry{Concentration: AI in Healthcare (Research) \; $\cdot$ \; Grade point average: 3.94 / 4.00}

\vspace{0.35em}
\entry{Federal University of Technology, Minna, Nigeria --- Bachelor of Engineering, Mechatronics Engineering}{Sep 2018 -- Jan 2025}
\subentry{First Class Honors \; $\cdot$ \; Cumulative grade point average: 4.60 / 5.00 \; $\cdot$ \; Concentration: Embedded Systems and Machine Learning}

\section{Research and Teaching Experience}

\entry{Research Intern --- Language and Spatial Laboratories}{June 2026 -- Present}
\subentry{Carnegie Mellon University $\cdot$ Supervisor: \href{https://scholar.google.com/citations?user=8PbgiPkAAAAJ&hl=en}{Prof. Prasenjit Mitra}}
\begin{bullets}
  \item Research on vision--language models for 3D medical imaging; lead developer of ICTC (see Research Projects).
  \item Built the full research pipeline, from data ingestion and the compression projector to multi-task training and evaluation across six visual token budgets.
\end{bullets}

\entry{Research Assistant (Intern)}{Oct 2023 -- Jan 2025}
\subentry{Artificial Intelligence for Clean Energy Limited, Minna, Nigeria $\cdot$ Supervisor: \href{https://scholar.google.com/citations?user=EPv13M8AAAAJ&hl=en}{Dr. Jibril Abdullahi Bala}}
\begin{bullets}
  \item Built a Flask and Random Forest energy-monitoring platform on Google Cloud (99\% uptime), improving monitoring accuracy by 30\%.
  \item Piloted drones for field data collection, cutting collection time by 30\% and improving data quality by 12\%.
\end{bullets}

\entry{Teaching Assistant --- AI System Design (04-652)}{Fall 2026}
\subentry{Carnegie Mellon University Africa $\cdot$ Supervisor: \href{https://scholar.google.com/citations?user=l_KmkVcAAAAJ&hl=en}{Prof. Hatem Haddad}}
\begin{bullets}
  \item Supporting a course on edge-to-cloud AI systems, from embedded machine learning on Raspberry Pi and Arduino to cloud-based model analysis.
\end{bullets}

\entry{Teaching Assistant --- Cloud Computing (15-319 / 15-619)}{Spring 2026}
\subentry{Carnegie Mellon University $\cdot$ Supervisor: \href{https://www.cs.cmu.edu/~msakr/}{Prof. Majd Sakr}}
\begin{bullets}
  \item Supported more than 200 undergraduate, master's and doctoral students through office hours, course-overview sessions and grading.
\end{bullets}

\entry{Peer Instructor (Volunteer) --- Introduction to Machine Learning}{Spring 2026}
\subentry{Carnegie Mellon University}
\begin{bullets}
  \item Led small-group tutoring on regression, classification, model evaluation and the underlying linear algebra and probability.
\end{bullets}

\section{Research Projects}

\entry{\href{https://github.com/Tech-sam-90/3d_compression}{ICTC: Instruction-Conditioned Token Compression for 3D CT Vision--Language Models}}{June 2026 -- Present}
\subentry{Language and Spatial Laboratories, Carnegie Mellon University $\cdot$ Advisor: \href{https://scholar.google.com/citations?user=8PbgiPkAAAAJ&hl=en}{Prof. Prasenjit Mitra}}
\stack{Python, PyTorch, Vision Transformers, Perceiver, LoRA, Hugging Face, CT-RATE}
\begin{bullets}
  \item Designed a two-stage token compressor for 3D chest CT report generation. A Perceiver stage condenses each slice to 32 latents, then an instruction-conditioned stage condenses the slices into a fixed budget of visual tokens for a LoRA-adapted language model.
  \item On the CT-RATE validation set (3,039 scans): ROUGE-L 0.273, BLEU-4 0.128 and RaTEScore 0.801 for full report generation.
  \item Built a multi-task instruction pipeline (report generation, entity questions, abnormality question answering) with a 368-test suite, evaluated across six token budgets from 16 to 512.
  \item Now adding instruction-conditioned plane selection, where the clinical question chooses an interpretable, anatomy-aware 3D viewing plane.
\end{bullets}

\entry{\href{https://github.com/Tech-sam-90/fpn-inceptentionnet}{FPN-Mamba: Linear-Time Bidirectional State Space Feature Pyramids for Small-Lesion Brain MRI Classification}}{Feb 2026 -- Present}
\subentry{Advisor: \href{https://scholar.google.com/citations?user=TEMlkDMAAAAJ&hl=en}{Prof. Moise Busogi}}
\stack{Python, PyTorch, Feature Pyramid Networks, Mamba state-space models, scikit-learn}
\begin{bullets}
  \item Combined a feature pyramid network with bidirectional Mamba state-space blocks to detect small brain metastases in MRI, in linear time.
  \item Against a re-implemented InceptentionNet (5 seeds $\times$ 3 folds), reaches AUC 0.982 vs 0.897 and smallest-lesion recall 0.870 vs 0.788 (all p \textless{} 0.02), and trains 3.7$\times$ faster per step than a matched self-attention design at 512 pixels.
  \item Accepted at the Africa in AI Workshop, NeurIPS 2026 --- \href{https://openreview.net/forum?id=te52LwNBSf}{OpenReview}.
\end{bullets}

\entry{\href{https://github.com/Tech-sam-90/malaria_fp_analysis}{Investigating White Blood Cells as a Source of False-Positive Malaria Parasite Detection in African Blood-Smear Images}}{Jan 2026 -- Aug 2026}
\stack{Python, PyTorch, Ultralytics YOLOv12, OpenCV, SciPy, DBSCAN, Weights \& Biases}
\begin{bullets}
  \item Ran a controlled study on 8,000 Giemsa-stained thick blood-smear images from Uganda and Ghana, comparing detectors trained with and without white blood cell labels.
  \item Seven independent spatial and statistical analyses (including Ripley's Cross-K, DBSCAN and Mann--Whitney U) refuted the common belief that white blood cells cause false-positive parasite detections, pointing to staining artefacts instead.
  \item Accepted for oral presentation at the MIRASOL Workshop, MICCAI 2026 --- \href{https://papers.miccai.org/miccai-2026-sat/MIRASOL_018.html}{paper} $\cdot$ \href{https://arxiv.org/abs/2609.29663}{arXiv:2609.29663}.
\end{bullets}

\entry{\href{https://github.com/Tech-sam-90/condition-aware_risk_CDSS}{Condition-Aware Intensive Care Unit Early-Warning System}}{Jan 2026 -- May 2026}
\stack{Python, PyTorch, TensorFlow, MIMIC-IV, FastAPI, Docker, WebSockets, scikit-learn}
\begin{bullets}
  \item Benchmarked six diagnosis-conditioned sequence models (LSTM, GRU, TCN, BiLSTM with attention, Transformer, TFT-style) for deterioration risk on MIMIC-IV intensive care data; the best, a BiLSTM with attention, reaches ROC-AUC 0.76.
  \item Deployed it as a Dockerized FastAPI service with WebSocket streaming inference, built for reproducibility under PhysioNet data governance.
\end{bullets}

\section{Publications}
\begin{enumerate}[leftmargin=1.4em, itemsep=0.25em, topsep=0.2em]
  \small
  \item \textbf{S.~A.~Adeniji}, G.~C.~Obasi, C.-V.~Ntwali, A.~M.~Iorumbur, C.~Raymond, L.~Uwimana, A.~T.~Issah. ``Investigating White Blood Cells as a Source of False-Positive Malaria Parasite Detection in African Blood-Smear Images.'' \textit{Medical Image Computing in Resource Constrained Settings Workshop \& Knowledge Interchange (MIRASOL), International Conference on Medical Image Computing and Computer Assisted Intervention (MICCAI) 2026}, Lecture Notes in Computer Science, vol. 17264, Springer (oral presentation). \href{https://papers.miccai.org/miccai-2026-sat/MIRASOL_018.html}{[MICCAI]} $\cdot$ \href{https://arxiv.org/abs/2609.29663}{[arXiv]} $\cdot$ \href{https://openreview.net/forum?id=DAv0RwrdIk}{[OpenReview]}
  \item S.~D.~Adetoro, \textbf{S.~A.~Adeniji}, M.~Busogi. ``FPN-Mamba: Linear-Time Bidirectional State Space Feature Pyramids for Small-Lesion Brain MRI Classification.'' \textit{Africa in AI Workshop, Conference on Neural Information Processing Systems (NeurIPS) 2026}. \href{https://openreview.net/forum?id=te52LwNBSf}{[OpenReview]}
  \item K.~E.~Jack, \textbf{S.~A.~Adeniji}. ``Convergence of Blockchain, Smart Microgrid and Energy Market.'' \textit{Electrica}, 24(2), 367--374, 2024. DOI: 10.5152/electrica.2024.23132.
  \item K.~E.~Jack, \textbf{S.~A.~Adeniji}. ``Development of an Artificial Intelligence-Based Infant's Emotion Recognition System for a Caregiving Robot.'' DOI: 10.5455/NJEAS.209325.
\end{enumerate}

\section{Academic Service and Research Training}
\begin{bullets}
  \item \textbf{SPRINT / SPARK AI for African Medical Imaging Knowledge Translation}, 2026 Cohort --- state-of-the-art methods for glioma brain-tumour segmentation and translation of imaging AI into African clinical practice; led a capstone on malaria parasite detection.
\end{bullets}

\section{Technical Skills}
\begin{bullets}
  \item \textbf{Languages and Platforms:} Python, Linux/Bash, Git, \LaTeX{}, Microsoft Azure, Google Cloud Platform.
  \item \textbf{Deep Learning:} PyTorch, TensorFlow, Vision Transformers, Mamba / state-space models, CNNs and Feature Pyramid Networks, vision--language models, LoRA / parameter-efficient fine-tuning, HuggingFace.
  \item \textbf{Medical Imaging and Data:} 3D CT and MRI pipelines, CT-RATE, MIMIC-IV, MONAI-style preprocessing, OpenCV, YOLO, scikit-learn, Weights \& Biases, Docker.
  \item \textbf{Foundations:} Linear algebra, probability and statistics, mathematical modelling, experimental design.
\end{bullets}

\section{Awards and Scholarships}
\begin{bullets}
  \item \textbf{Mastercard Foundation Scholarship}, Carnegie Mellon University (2025), for the master's in Engineering Artificial Intelligence.
  \item \textbf{MINDS Scholarship} (Mandela Institute of Developmental Studies), Carnegie Mellon University (2026).
  \item \textbf{Best Graduating Student Project}, Department of Mechatronics Engineering, Federal University of Technology, Minna (2024): a nanny robot for childcare and monitoring.
\end{bullets}

\end{document}
